% Part: proof-theory
% Chapter: proof-search
% Section: tableaux

\documentclass[../../../include/open-logic-section]{subfiles}

\begin{document}

\olfileid{pt}{ps}{tab}

\olsection{ট্যাবলো}

ট্যাবলো পদ্ধতিগুলি সিকোয়েন্ট কলনের ঘনিষ্ঠ !!{proof} পদ্ধতি;
হাতে বা সফটওয়্যারে প্রমাণ-অন্বেষণের পক্ষে এগুলি বিশেষ উপযোগী।
সিকোয়েন্টের বৃক্ষের বদলে ট্যাবলোতে !!a{proof} হলো
!!{formula}s-এর বৃক্ষ। তবে স্বাভাবিক নিষ্পাদনের বিপরীতে
এর বিধিগুলি সিকোয়েন্ট কলনের বিধির প্রতিরূপ। ট্যাবলো
!!{proof} মূলত মডেল খোঁজার একটি স্পষ্ট অনুসন্ধান। তাই একে
কখনো \emph{অর্থতাত্ত্বিক} ট্যাবলো, কখনো \emph{সত্য-বৃক্ষ} বলে।

একটি \emph{চিহ্নিত} ট্যাবলো হলো চিহ্নিত !!{formula}s-এর
বৃক্ষ; চিহ্নিত সূত্র $\sFmla{\True}{!A}$ বা
$\sFmla{\False}{!A}$ আকারের হয়। বৃক্ষ নিচের দিকে বাড়ে।
শুরুতে থাকে কী !!{prove} করতে চাই তা নির্দিষ্ট করা
!!{formula}s। যেমন, $!A, !B \Sequent !C, !D$
সিকোয়েন্টটি !!{prove} করতে শুরুতে লিখি
$\sFmla{\True}{!A}, \sFmla{\True}{!B},
\sFmla{\False}{!C}, \sFmla{\False}{!D}$।
$!A$ ও $!B$ সত্য, $!C$ ও $!D$ মিথ্যা করে এমন
!!^a{structure} হলো $!A, !B \Sequent !C, !D$ অবৈধ
দেখানো !!a{structure}।
% BN-SRC-602 TARGET-CORRECTION-BEGIN
\par\smallskip\noindent\textnormal{\small\textbf{উৎস-সংশোধন (BN-SRC-602):} ডান পাশের $!D$-কে মিথ্যা দেখানোর জন্য শুরুর চিহ্ন $\False$; উৎসে ভুল করে $\True$ ছিল।}
% BN-SRC-602 TARGET-CORRECTION-END

ট্যাবলোর পাতার (সবচেয়ে নিচের) গিঁট শাখা-প্রসারণ বিধি
দিয়ে বাড়ানো যায়। বিধি একই শাখায় আরও গিঁট যোগ করতে
পারে, কিংবা নিচে দুটি সহোদর গিঁট যোগ করে শাখাটিকে
দুই শাখায় ভাগ করতে পারে। সেই শাখায় পাতার উপরে থাকা
যেকোনো চিহ্নিত !!{formula}-য় বিধি প্রয়োগ করা যায়।
ধ্রুপদি ট্যাবলো পদ্ধতি \Log{Tc}-র শাখা-প্রসারণ বিধি
\olref{tab:Tc}-এ দেওয়া আছে।

\subfile{rules-Tc}

দেখা যায়, বিধিগুলি \Log{G3c}-র বিধি উল্টো দিকে
প্রয়োগেরই রূপ। $\True$ ও $\False$ চিহ্ন যথাক্রমে
সিকোয়েন্টের বাঁ ও ডান দিকে কোন !!a{formula} প্রধান,
তা বোঝায়।

ট্যাবলোর একটি শাখায় $\sFmla{\True}{!A}$ ও
$\sFmla{\False}{!A}$ দুটিই থাকলে, অথবা $\sFmla{\True}{\lfalse}$
থাকলে শাখাটি \emph{বদ্ধ}। প্রতিটি শাখা বদ্ধ হলে গোটা ট্যাবলো
বদ্ধ। যেমন,
$!A \lor (!B \land !C) \Sequent (!A \lor !B) \land
(!A \lor !C)$-র একটি বদ্ধ ট্যাবলো নিচে আছে।
এর শুরুতে $\sFmla{\True}{!A \lor (!B \land
!C)}$ ও $\sFmla{\False}{(!A \lor !B) \land (!A \lor !C)}$:
% BN-SRC-603 TARGET-CORRECTION-BEGIN
\par\smallskip\noindent\textnormal{\small\textbf{উৎস-সংশোধন (BN-SRC-603):} উদাহরণটির সব শাখা বদ্ধ এবং ঠিক আগেই বদ্ধ ট্যাবলো সংজ্ঞায়িত; উৎসের “clause tableau” এখানে “closed tableau” অর্থে পড়া হয়েছে।}
% BN-SRC-603 TARGET-CORRECTION-END
% BN-SRC-889 TARGET-CORRECTION-BEGIN
\par\smallskip\noindent\textnormal{\small\textbf{পাঠ-সংশোধন (BN-SRC-889):} মিথ্যা-সূত্রকে সত্য দাবি করা শাখাও বদ্ধ; সিকোয়েন্ট কলনে এটি মিথ্যা-সূত্রের বাঁ স্বতঃসিদ্ধের ক্ষেত্র। উৎসের শুধু বিপরীত-চিহ্নের জোড়ার শর্তে এই ক্ষেত্র বাদ ছিল।}
% BN-SRC-889 TARGET-CORRECTION-END
\begin{oltableau}
  [\sFmla{\True}{\formula{A} \lor (\formula{B} \land \formula{C})},
    just=\TAss, checked
    [\sFmla{\False}{(\formula{A} \lor \formula{B}) \land
        (\formula{A} \lor \formula{C})}, just=\TAss, checked
      [\sFmla{\True}{\formula{A}}, just={\TRule{\True}{\lor}[1]}
        [\sFmla{\False}{\formula{A} \lor \formula{B}},
          just={\TRule{\False}{\land}[2]}, checked
          [\sFmla{\False}{\formula{A}},
            just={\TRule{\False}{\lor}[4]}
            [\sFmla{\False}{\formula{B}},
              just={\TRule{\False}{\lor}[4]}, close]
          ]
        ]
        [\sFmla{\False}{\formula{A} \lor \formula{C}},
          just={\TRule{\False}{\land}[2]}, checked
          [\sFmla{\False}{\formula{A}},
            just={\TRule{\False}{\lor}[4]}
            [\sFmla{\False}{\formula{C}},
              just={\TRule{\False}{\lor}[4]}, close]
          ]
        ]
      ]
      [\sFmla{\True}{\formula{B} \land \formula{C}},
        just={\TRule{\True}{\lor}[1]}, checked
        [\sFmla{\False}{\formula{A} \lor \formula{B}},
          just={\TRule{\False}{\land}[2]}, checked
          [\sFmla{\True}{\formula{B}},
            just={\TRule{\True}{\land}[3]}, move by=2
            [\sFmla{\True}{\formula{C}},
              just={\TRule{\True}{\land}[3]}
              [\sFmla{\False}{\formula{A}},
                just={\TRule{\False}{\lor}[4]}
                [\sFmla{\False}{\formula{B}},
                  just={\TRule{\False}{\lor}[4]},close
                ]
              ]
            ]
          ]
        ]
        [\sFmla{\False}{\formula{A} \lor \formula{C}},
          just={\TRule{\False}{\land}[2]}, checked
          [\sFmla{\True}{\formula{B}},
            just={\TRule{\True}{\land}[3]}, move by=2
              [\sFmla{\True}{\formula{C}},
                just={\TRule{\True}{\land}[3]}
                [\sFmla{\False}{\formula{A}},
                  just={\TRule{\False}{\lor}[4]}
                  [\sFmla{\False}{\formula{C}},
                  just={\TRule{\False}{\lor}[4]},close
                ]
              ]
            ]
          ]
        ]
      ]
    ]
  ]
\end{oltableau}

টিকচিহ্ন বোঝায় যে নিচের সব শাখায় সংশ্লিষ্ট
!!a{formula}-য় বিধিটি প্রয়োগ করা হয়েছে। $\otimes$
চিহ্নটি শাখা বদ্ধ হওয়া বোঝায়। বিধির পরে থাকা
লাইন-সংখ্যা নির্দেশ করে, কোন চিহ্নিত !!{formula}s-এ
বিধিটি প্রয়োগ হয়েছে—অর্থাৎ দেখানো লাইনে একই
শাখায় থাকা !!{formula}-য়।

ট্যাবলোতে প্রমাণ-অন্বেষণ সহজ। পর্যায়ক্রমে ট্যাবলো
তৈরি করি। $0$-তম পর্যায়ে শুরুর !!{formula}s উপরে
লিখি। $n+1$-তম পর্যায়ে প্রতিটি পাতার গিঁটে
পরের অনুচ্ছেদে বিশেষভাবে সামলানো পরিমাণায়ক সূত্রগুলি বাদ দিয়ে,
এখনও টিকচিহ্ন না দেওয়া প্রথম (সবচেয়ে উপরের)
প্রয়োগযোগ্য অ-পারমাণবিক চিহ্নিত সূত্রটি খুঁজি। এমন সূত্র থাকলে
সেটিতে শাখা-প্রসারণ বিধি প্রয়োগ করি। আলোচ্য চিহ্নিত সূত্রটি থাকা প্রতিটি
শাখায় বিধি প্রয়োগ হয়ে গেলে তাতে টিকচিহ্ন দিই।
(উপরের উদাহরণটি এই অ্যালগরিদমে তৈরি।)
% BN-SRC-892 TARGET-CORRECTION-BEGIN
\par\smallskip\noindent\textnormal{\small\textbf{পাঠ-সংশোধন (BN-SRC-892):} পারমাণবিক সূত্রের শাখা-প্রসারণ বিধি নেই। তাই চিহ্ন না-পড়া প্রথম সূত্রের বদলে প্রথম প্রয়োগযোগ্য সূত্র বাছতে হয়; এমন সূত্র না থাকলে সাধারণ প্রসারণের ধাপ বাদ যায়।}
% BN-SRC-892 TARGET-CORRECTION-END

$\sFmla{\True}{\lforall}$ ও
$\sFmla{\False}{\lexists}$ আকারের চিহ্নিত
!!{formula}s-এ কখনো টিকচিহ্ন দেওয়া হয় না;
তাই এগুলিকে বিশেষভাবে সামলাতে হয়। এগুলি
থাকলে নিশ্চিত করতে হবে, সব প্রযোজ্য পদের জন্য
$\TRule{\True}{\forall}$ ও
$\TRule{\False}{\lexists}$ বিধি একসময় প্রয়োগ
হচ্ছে। যেমন, প্রতি পর্যায়ে প্রতিটি শাখার
এমন সব সূত্রে বিধি প্রয়োগ করলেই চলে (এই
আকারের নয় এমন সর্বোচ্চ চিহ্নিত !!{formula}
সামলানোর পরে)। ব্যবহারের পদ~$t$ হলো
শাখায় ইতিমধ্যে থাকা !!{constant}s (কোনো
!!{constant} না থাকলে~$\Obj c_0$)
ও শুরুর !!{formula}s-এ থাকা !!{function}s
দিয়ে গড়া প্রথম পদ~$t$, যার সংশ্লিষ্ট চিহ্নিত
!!{formula} $\sFmla{\True}{!B(t)}$
বা~$\sFmla{\False}{!B(t)}$ ওই শাখায়
এখনও নেই। এমন পদ থাকলেই এই বিধি প্রয়োগ করি।
% BN-SRC-890 TARGET-CORRECTION-BEGIN
\par\smallskip\noindent\textnormal{\small\textbf{পাঠ-সংশোধন (BN-SRC-890):} নতুন উদাহরণ-সূত্র আর না থাকলে ওই ধাপে পরিমাণায়ক বিধি প্রয়োগ হয় না। পুনরাবৃত্ত পরিমাণায়ক সূত্রে কখনো টিকচিহ্ন পড়ে না বলে সসীম খোলা শাখার শেষ অবস্থা সব সূত্রে টিকচিহ্ন দিয়ে সংজ্ঞায়িত করা যায় না; সব প্রযোজ্য প্রসারণ সম্পন্ন থাকা প্রয়োজন।}
% BN-SRC-890 TARGET-CORRECTION-END

এই প্রমাণ-অন্বেষণে বদ্ধ ট্যাবলো না মিললে,
হয় !!{formula}s থেকে কোনো বিধিতে আর নতুন গিঁট যোগ করা যায় না
এমন একটি সসীম খোলা শাখা থাকে,
নয় অনুসন্ধানে একটি অনন্ত, সসীম-শাখাবিশিষ্ট
বৃক্ষ তৈরি হয়; তাতে অবশ্যই অনন্ত শাখা আছে।
\olref[cpl]{sec}-এর মতোই এমন
!!a{structure}~$\Struct{M}$ সংজ্ঞায়িত
করা যায় যাতে শাখায় $\sFmla{\True}{!A}$
থাকলে $\Sat{M}{!A}$ এবং
$\sFmla{\False}{!A}$ থাকলে
$\Sat/{M}{!A}$। (ধরি, $\Theta =
\Setabs{!A}{\sFmla{\True}{!A} \text{ শাখায় আছে}}$ এবং $\Xi =
\Setabs{!A}{\sFmla{\False}{!A} \text{ শাখায় আছে}}$।)
% BN-NORM-175: Bengali display prose retained; optional reader annotation withdrawn.

এই প্রমাণ-অন্বেষণে তৈরি ট্যাবলোর প্রতিটি
গিঁটে একটি সিকোয়েন্ট বসিয়ে সহজেই
\Log{G3c} প্রমাণে অনুবাদ করা যায়।
শুরুর সূত্রগুলি $\Gamma \Sequent
\Delta$ সিকোয়েন্টের হলে, প্রতিটি শুরুর
সূত্রে $\Gamma \Sequent
\Delta$ লেবেল দিই। চিহ্নিত
!!{formula} $\sFmla{\True}{!A}$-তে
শাখা-প্রসারণ বিধি প্রয়োগ করে শাখা বাড়ালে
পাতার লেবেল $!A, \Pi \Sequent \Lambda$;
চিহ্নিত !!{formula}
$\sFmla{\False}{!A}$-তে প্রয়োগ করলে
পাতার লেবেল $\Pi \Sequent
\Lambda, !A$। ট্যাবলোতে
\TRule{\True}{\star} বিধি প্রয়োগের
স্থানে সিকোয়েন্টে \LeftR{\star}
বিধি প্রয়োগ করি (প্রধান সূত্র~$!A$);
\TRule{\False}{\star}-র স্থানে
\RightR{\star} প্রয়োগ করি।
বিধির পূর্বধারণাগুলি হলো ট্যাবলোতে
যোগ হওয়া সংশ্লিষ্ট গিঁটগুলির
লেবেল-সিকোয়েন্ট। কোনো ট্যাবলো
$\sFmla{\True}{!A}$ ও
$\sFmla{\False}{!A}$ দিয়ে বদ্ধ হলে
শেষ পাতার লেবেল
$!A, \Pi \Sequent \Lambda,
!A$ আকারের সিকোয়েন্ট। মিথ্যা-সূত্রের জন্য শাখা বদ্ধ হলে
শেষ সিকোয়েন্ট মিথ্যা-সূত্রের বাঁ স্বতঃসিদ্ধ।
ট্যাবলোর পরপর কয়েকটি গিঁটে
একই সিকোয়েন্ট বসতে পারে;
পুনরাবৃত্তগুলি বাদ দিই।
যেমন, উপরের ট্যাবলো নিচের
!!{proof}-এ অনূদিত হয়:
\[\small
\Axiom$!A \fCenter !A, !B$
\RightLabel{\RightR{\lor}}
\UnaryInf$!A \fCenter !A \lor !B$
\Axiom$!A \fCenter !A, !C$
\RightLabel{\RightR{\lor}}
\UnaryInf$!A \fCenter !A \lor !C$
\RightLabel{\RightR{\land}}
\BinaryInf$!A \fCenter (!A \lor !B) \land (!A \lor !C)$
\Axiom$!B, !C \fCenter !A, !B$
\RightLabel{\RightR{\lor}}
\UnaryInf$!B, !C \fCenter !A \lor !B$
\RightLabel{\LeftR{\land}}
\UnaryInf$!B \land !C \fCenter !A \lor !B$
\Axiom$!B, !C \fCenter !A, !C$
\RightLabel{\RightR{\lor}}
\UnaryInf$!B, !C \fCenter !A \lor !C$
\RightLabel{\LeftR{\land}}
\UnaryInf$!B \land !C \fCenter !A \lor !C$
\RightLabel{\RightR{\land}}
\BinaryInf$!B \land !C \fCenter (!A \lor !B) \land
(!A \lor !C)$
\RightLabel{\LeftR{\lor}}
\BinaryInf$!A \lor (!B \land !C) \fCenter (!A \lor !B) \land
(!A \lor !C)$
\DisplayProof
\]

\end{document}
